% W4i27_P_updates.tex
% DFD 17-Agent Research Programme — Iteration 27
% Agents 6 (P1-P4 Dead Patents), 7 (P5-P7 ET/CE), 8 (P8-P11 SQMS/Dark SRF)
% Date: 2026-03-21

\documentclass[11pt,a4paper]{article}
\usepackage[utf8]{inputenc}
\usepackage{amsmath,amssymb,amsfonts}
\usepackage{hyperref}
\usepackage{booktabs}
\usepackage{longtable}
\usepackage{geometry}
\geometry{margin=2.5cm}

\title{DFD i27 Patent Updates\\Agents 6, 7, 8}
\author{DFD 17-Agent Research Programme}
\date{2026-03-21}

\begin{document}
\maketitle
\tableofcontents
\newpage

%=============================================================================
\section{Agent 6: Patents P1--P4 --- Dead Patents Review}
%=============================================================================

\subsection{Mandate}
Monitor developments relevant to dead patents P1 (Field Drag Cloaking), P2 (Resonators/Metamaterials), P3 (Quantum Control), P4 (Energy Coupling/Power). Assess whether any new physics could revive these patents.

\subsection{New Literature}

\begin{enumerate}

\item \textbf{Fermilab News Release (2026-03-17)} --- ``Fermilab-led study advances development of sophisticated quantum sensors to track high-energy particles and detect dark matter.'' Reports on transmon-based quantum sensors for particle tracking and dark matter detection using SRF cavities. Quantum-enhanced sensing with single microwave photon counters. \textbf{Grade: C for P3.} Quantum control via SRF+transmon systems is advancing rapidly. However, P3's specific claims about DFD-mediated quantum control remain unsupported by data. The technology is relevant to P9/P11 (Agent 8) rather than P3.

\item \textbf{Zhai et al.\ (2026), PRB accepted} --- ``Quantum-enhanced sensing of axion dark matter with a transmon-based single microwave photon counter.'' Demonstrates sub-quantum-limit sensitivity in SRF cavities using transmon qubits. \textbf{Grade: C for P3/P4.} The transmon-SRF interface is maturing. If DFD's $\psi$-field couples to EM fields in SRF cavities (P11), transmon readout could probe it. But this is speculative for P3/P4's specific claims.

\item \textbf{Differential torsion sensor (2025), arXiv:2402.08935, PRD 111, 063064} --- ``Differential torsion sensor for direct detection of ultralight vector dark matter.'' New sensor design for ultralight DM detection at sub-neV masses. \textbf{Grade: C for P1.} While not directly relevant to cloaking, the ultralight field detection technology could in principle probe DFD's $\psi$-field gradient if it exists at terrestrial scales. P1 remains DEAD: no mechanism to generate controllable $\psi$-field gradients.

\end{enumerate}

\subsection{Patent Status Assessment}

\begin{center}
\begin{tabular}{lccc}
\toprule
\textbf{Patent} & \textbf{Status} & \textbf{Change from i26} & \textbf{Revival Prospect} \\
\midrule
P1 (Field Drag Cloaking) & DEAD & No change & None \\
P2 (Resonators/Metamaterials) & DEAD & No change & None \\
P3 (Quantum Control) & DEAD & No change & Remote (needs P11 first) \\
P4 (Energy Coupling/Power) & DEAD & No change & None \\
\bottomrule
\end{tabular}
\end{center}

\textbf{Assessment:} No new developments that would revive any of P1--P4. The quantum sensor advances (transmon+SRF) are relevant to ALIVE patents (P9, P11) but do not address the fundamental problems with P1--P4: these patents assume controllable macroscopic $\psi$-field manipulation, which remains physically implausible. \textbf{All four patents remain DEAD. Grade: N/A (monitoring only).}

%=============================================================================
\section{Agent 7: Patents P5--P7 --- ET/CE Timeline and GW Detectors}
%=============================================================================

\subsection{Mandate}
Track Einstein Telescope and Cosmic Explorer timelines. Assess P5 (Geometry/Nav/Timing), P6 (Quantum Sensors/Imaging), P7 (Energy Harvesting/Storage) relevance. Focus on third-generation GW detector capabilities.

\subsection{New Literature}

\begin{enumerate}

\item \textbf{Einstein Telescope Consortium (2025), Science Case Update} --- ET site decision expected 2026 (Sardinia vs.\ Euregio Meuse-Rhine). Construction start 2028. First light 2035. Budget: 1.8 billion EUR. Three nested detectors, 10 km arms, cryogenic optics (10--20 K). Sensitivity: factor 10 improvement over Advanced LIGO/Virgo. Low-frequency cutoff: 3 Hz (vs.\ 10 Hz for current detectors). \textbf{Grade: B for P5.} ET's improved low-frequency sensitivity will enable better measurement of tidal deformability in NS mergers, relevant to DFD's predictions for neutron star structure via the scalar field. P5's navigation/timing claims remain NICHE: ET does not measure position or timing with the precision needed for P5's GPS-alternative concept.

\item \textbf{Cosmic Explorer Project Status (2025), CE-P2500001} --- NSF awarded \$9M for conceptual design. Two interferometers: 40 km and 20 km arms. Operational target: mid-2030s to 2040. Sensitivity: 10$\times$ LIGO. Will observe GWs from the edge of the observable universe. International partners: UK, Germany, Australia, Canada. \textbf{Grade: B for P5/P7.} CE's factor-10 sensitivity improvement is relevant to DFD prediction P31 (gravitational slip): combined ET+CE observations of binary neutron star mergers at multiple redshifts could constrain the scalar field's effect on GW propagation. P7 (energy harvesting) remains NICHE: no mechanism to extract energy from GW detectors.

\item \textbf{Verma et al.\ (2025), arXiv:2509.05027} --- ``Unique gravitational wave signatures of GLPV scalar-tensor theories.'' Discovers a new scalar-scalar-tensor interaction unique to beyond-Horndeski (GLPV) theories. Predicts induced GW spectral density $\propto f^5$ for scale-invariant primordial spectrum. Higher-derivative terms enhance GW production. \textbf{Grade: A for P5.} \emph{New DFD implication}: although DFD sits within Horndeski (not GLPV), the disformal sector contributes to scalar-induced GWs. If DFD's disformal factor $\Sigma(y)$ generates a characteristic induced-GW spectrum, ET/CE could detect it. The $f^5$ scaling of GLPV provides a discriminator: DFD (Horndeski) would predict a \emph{different} scaling. This is a \textbf{new testable distinction}.

\item \textbf{ET October 2025 Status Report} --- European governments (Netherlands, Belgium, Germany, Italy) have expressed national-level support. Site announcement imminent (2026). \textbf{Grade: B (timeline).} P5 and related patents that depend on third-generation GW science have a 2035+ timeline.

\end{enumerate}

\subsection{Patent Status Assessment}

\begin{center}
\begin{tabular}{lccc}
\toprule
\textbf{Patent} & \textbf{Status} & \textbf{Change from i26} & \textbf{Timeline} \\
\midrule
P5 (Geometry/Nav/Timing) & ALIVE & No change & ET 2035+, CE 2040+ \\
P6 (Quantum Sensors) & NICHE & No change & Depends on P11 \\
P7 (Energy Harvesting) & ALIVE & No change & ET/CE era \\
\bottomrule
\end{tabular}
\end{center}

\textbf{Key Finding: GLPV-vs-Horndeski GW signature.} Verma et al.\ (2509.05027) identify a unique beyond-Horndeski GW signature ($f^5$ scaling). DFD, being Horndeski-class, would predict a different scaling. This provides a \emph{new discriminator} between DFD and competing beyond-Horndeski theories, testable with ET/CE. \textbf{New sub-prediction to add to P5.}

\textbf{Status: P5 and P7 remain ALIVE. ET site decision 2026, construction 2028, first light 2035. CE mid-2030s--2040. DFD GW predictions have a 10+ year timeline but are well-defined. Grade B.}

%=============================================================================
\section{Agent 8: Patents P8--P11 --- SQMS, Dark SRF, and Transmon Probes}
%=============================================================================

\subsection{Mandate}
Track SQMS Center developments, Dark SRF experimental results, and transmon-SRF coupling measurements. Assess P8 (Field Communications), P9 (Field Navigation), P10 (Field Computing), P11 (EM Coupling/SRF).

\subsection{New Literature}

\begin{enumerate}

\item \textbf{Kalia et al.\ (2026), PRL 136, 111801, arXiv:2510.02427} --- ``Improved Dark Photon Sensitivity from the Dark SRF Experiment.'' Refined dark-photon exclusion from pathfinder run. Order-of-magnitude improvement over previously reported bounds. World-leading constraint on non-dark-matter dark photons below 6 $\mu$eV. Photon mass: $m_\gamma < 2.9 \times 10^{-48}$ g (best laboratory limit). Key innovation: theoretical modeling of frequency instability in high-Q resonant experiments. \textbf{Grade: A for P9/P11.} The frequency-instability modeling (Kalia+ 2504.15307) is \emph{directly transferable} to DFD P9's Q-ratio measurement. If $\psi$-field coupling exists, it would manifest as anomalous frequency instability in SRF cavities. The order-of-magnitude improvement in sensitivity brings P9's predicted Q-ratio anomaly ($Q_\mathrm{obs}/Q_\mathrm{BCS} = 0.52 \pm 0.08$ vs.\ 3.60 $\pm$ 0.10) closer to testability.

\item \textbf{Kalia et al.\ (2025), arXiv:2504.15307, PRD} --- ``Modeling frequency instability in high-quality resonant experiments.'' Develops theoretical framework for frequency jittering in SRF cavities. Microscopic deformations cause stochastic frequency shifts with amplitude 20$\times$ linewidth. Timescale of jittering determines effect on power accumulation. Constraint improvement: 4 orders of magnitude in signal-to-noise. \textbf{Grade: A for P9/P11.} This is the \emph{noise model} that P9 needs. To detect DFD's predicted Q-ratio anomaly, one must first model and subtract all known sources of frequency instability. Kalia et al.\ provide this model. The residual, after subtraction, would be the DFD signal.

\item \textbf{Zhai et al.\ (2026), PRB accepted} --- ``Quantum-enhanced sensing of axion dark matter with a transmon-based single microwave photon counter.'' Achieves sub-quantum-limit sensitivity. Demonstrates transmon readout of SRF cavity photons at single-quantum level. \textbf{Grade: A for P11.} DFD's transmon-$\psi$ probe concept (proposed i24): couple a transmon qubit to an SRF cavity operating at the DFD-predicted anomalous Q-ratio. If the $\psi$-field shifts the cavity resonance, the transmon would detect single-photon-level anomalies. Zhai et al.\ demonstrate that this readout sensitivity is \emph{already achieved}.

\item \textbf{SQMS Center News (2026-03)} --- Fermilab leads project to develop novel quantum sensors for high-energy particles and dark matter using SRF+transmon technology. Extended SQMS roadmap includes dark photon searches and fundamental physics measurements. \textbf{Grade: B for P9/P11.}

\item \textbf{OSTI Report (2025)} --- ``A Flux-Tunable cavity for Dark matter detection.'' Describes tunable SRF cavity for scanning dark photon mass range. Flux-tuning allows frequency adjustment without mechanical changes. \textbf{Grade: B for P9.} Flux-tunable cavities could scan for DFD's predicted Q-ratio anomaly across multiple frequencies, testing whether the anomaly is frequency-dependent (as DFD predicts: the $\psi$-coupling should scale with cavity mode structure).

\end{enumerate}

\subsection{Analysis: P9 Q-Ratio Measurement Feasibility}

The convergence of three developments makes P9 experimentally closer than ever:

\begin{enumerate}
\item \textbf{Noise model:} Kalia et al.\ (2504.15307) provide the frequency-instability framework.
\item \textbf{Sensitivity:} Dark SRF (2510.02427) achieves order-of-magnitude improvement.
\item \textbf{Readout:} Zhai et al.\ demonstrate single-photon transmon readout.
\end{enumerate}

DFD predicts $Q_\mathrm{obs}/Q_\mathrm{BCS} = 0.52 \pm 0.08$ (anomalous suppression from $\psi$-field coupling to superconducting condensate). BCS theory predicts $Q_\mathrm{obs}/Q_\mathrm{BCS} = 3.60 \pm 0.10$. The $>$30$\sigma$ discrepancy is an extraordinary claim requiring extraordinary evidence. The experimental path:

\begin{itemize}
\item Operate two identical SRF cavities: one in a magnetically shielded environment, one exposed to varying gravitational potential (e.g., different heights, proximity to massive objects).
\item Model frequency instability using Kalia framework.
\item Use transmon readout to detect any differential Q-ratio shift.
\item If $\Delta Q/Q \neq 0$ correlates with gravitational potential difference, this is a DFD signature.
\end{itemize}

\textbf{Timeline:} Current SQMS infrastructure could support a dedicated P9 experiment within 2--3 years if prioritised. The technology is ready; the bottleneck is allocation of beamtime and personnel.

\subsection{Patent Status Assessment}

\begin{center}
\begin{tabular}{lccl}
\toprule
\textbf{Patent} & \textbf{Status} & \textbf{Change from i26} & \textbf{Notes} \\
\midrule
P8 (Field Comms) & NICHE & No change & Requires P11 validation \\
P9 (Field Nav) & ALIVE & Strengthened & Dark SRF noise model + sensitivity \\
P10 (Field Computing) & NICHE & No change & Requires P11 validation \\
P11 (EM Coupling/SRF) & ALIVE & Strengthened & Transmon readout achieved \\
\bottomrule
\end{tabular}
\end{center}

\textbf{Key Finding: The three technological advances (noise modeling, sensitivity improvement, single-photon readout) form a complete experimental toolkit for testing P9/P11. The SQMS Center has all the necessary infrastructure. A dedicated DFD-targeted experiment is now feasible.}

\textbf{Status: P9 and P11 both STRENGTHENED. Dark SRF PRL provides the sensitivity; frequency-instability model provides the noise framework; transmon readout provides the detector. Grade A.}

%=============================================================================
\section{Summary}
%=============================================================================

\begin{center}
\begin{tabular}{lccc}
\toprule
\textbf{Agent} & \textbf{Grade} & \textbf{New Papers} & \textbf{Key Finding} \\
\midrule
6 (P1--P4) & N/A & 3 & All dead; no revival \\
7 (P5--P7) & B & 4 & GLPV vs Horndeski GW discriminator \\
8 (P8--P11) & A & 5 & Complete P9/P11 experimental toolkit \\
\bottomrule
\end{tabular}
\end{center}

\end{document}
